% ===========================================================================
% Section 3: System Design (~2 pages)
% ===========================================================================

% --- Architecture overview figure ---
% TODO: Create architecture diagram
% \begin{figure*}[t]
%   \centering
%   \includegraphics[width=\textwidth]{fig_architecture.pdf}
%   \caption{\system{} architecture overview. Phase~A constructs a
%   benchmark-grounded Knowledge Base. Phase~B trains ML classifiers on
%   gold and silver labels. Phase~C performs online inference: the ML
%   classifier detects bottleneck patterns, SHAP identifies key counters,
%   RAG retrieves relevant KB entries, and the LLM generates code-level
%   recommendations.}
%   \label{fig:architecture}
% \end{figure*}

% --- 3.1 Phase A: Benchmark Execution + KB Construction ---
% \subsection{Phase A: Knowledge Base Construction}
% \label{sec:design:kb}
%
% TODO: Describe benchmark execution (IOR, mdtest, DLIO)
% - Parameter sweeps to create known bottleneck patterns
% - Darshan log collection for each configuration
% - Knowledge Base: {Darshan signature <-> code cause <-> validated fix}
% - Verification: apply fix, re-run benchmark, measure improvement

% --- 3.2 Phase B: ML Classifier Training ---
% \subsection{Phase B: ML Classifier Training}
% \label{sec:design:ml}
%
% TODO: Describe training data construction
% - Gold labels: benchmark-verified (high quality, small quantity)
% - Silver labels: Drishti heuristic on production logs (lower quality, large quantity)
% - Cleanlab for noise handling in silver labels
% - Multi-label classification (each sample can have multiple bottleneck types)
% - Model comparison: XGBoost, LightGBM, RF, MLP

% --- 3.3 Phase C: Online Inference ---
% \subsection{Phase C: Online Inference}
% \label{sec:design:inference}
%
% TODO: Describe inference pipeline
% - Input: new Darshan log
% - Step 1: Feature extraction (186 features)
% - Step 2: ML classifier predicts bottleneck labels
% - Step 3: SHAP attribution identifies key counters
% - Step 4: RAG retrieves relevant KB entries using SHAP-weighted query
% - Step 5: LLM generates code-level recommendations grounded in KB
% - Output: bottleneck report + specific code fixes
